% Hodina 59 podle Rozpis_vyuky_MAT_1_rocnik.tex, řádek 59.
% Hlavní podklad: 1_rocnik/pripravy/Podklady_pro_tvorbu_priprav_Funkce.md.
% Odborný zdroj: ucebnice/02_Funkce_a_rovnice/Zaklady_funkci/Funkce.tex,
% části „Co je funkce“ a „Způsoby zadání funkce“.
% Metodické srovnání s lokálními PDF Realisticky.cz (čten obsah, ne přebírání úloh):
% 020103_Zobrazeni/020103_Lekce.pdf – konečná zobrazení a diagramy přiřazení;
% 020104_Funkce_definicni_obor/020104_Lekce.pdf – příklady různých vstupů,
% jednoznačnost, zadání předpisem/tabulkou/grafem a pedagogické poznámky;
% 020105_Graf_funkce_I/020105_Lekce.pdf – čtení dvojic a hodnot z grafu.
% Výklad, příklady, diagramy a grafy níže jsou autorským zpracováním.

\documentclass[11pt,a4paper]{article}
% Spolecny styl pro ucitelske pripravy MAT (Work/MAT)
% Verze 1.1, 2026-10-23
\NeedsTeXFormat{LaTeX2e}
\ProvidesPackage{priprava}[2026/10/23 v1.1 Sablona ucitelskych priprav MAT]

% Zakladni typografie a cestina
\RequirePackage[T1]{fontenc}
\RequirePackage[utf8]{inputenc}
\RequirePackage[czech]{babel}
\RequirePackage{lmodern}
\RequirePackage{microtype}
\RequirePackage{amsmath,amssymb}
\RequirePackage{xcolor}
\RequirePackage{enumitem}
\RequirePackage[most]{tcolorbox}
\RequirePackage{tabularx,booktabs,array}
\RequirePackage{fancyhdr}
\RequirePackage{needspace}
\RequirePackage{geometry}
\RequirePackage[hidelinks]{hyperref}

\geometry{a4paper,top=17mm,bottom=18mm,left=23mm,right=23mm}
\definecolor{notegray}{gray}{0.39}
\setlength{\parindent}{0pt}
\setlength{\parskip}{2.5pt plus 1pt minus 0.5pt}
\setlist[itemize]{topsep=2pt,itemsep=1pt,parsep=0pt,leftmargin=1.45em}
\setlist[enumerate]{topsep=2pt,itemsep=2pt,parsep=0pt,leftmargin=1.65em}
\setcounter{secnumdepth}{0}
\renewcommand{\arraystretch}{1.12}

% Zahlavi: vlevo rocnik | globalni tema, vpravo cislo: nazev hodiny
% Zapati zustava zcela prazdne.
\pagestyle{fancy}
\fancyhf{}
\renewcommand{\headrulewidth}{0.35pt}
\renewcommand{\footrulewidth}{0pt}
\setlength{\headheight}{14pt}
\newcommand{\nastavPripravu}[4]{%
  \fancyhead[L]{\small #1\enspace\textbar\enspace #2}%
  \fancyhead[R]{\small #3: #4}%
}

% Minimalisticke cernobile boxy. Kazda definice zvlast.
\tcbset{commonlesson/.style={enhanced,breakable=false,colback=white,colframe=black,
  colbacktitle=white,coltitle=black,boxrule=0.45pt,
  arc=0.8mm,outer arc=0.8mm,left=3mm,right=3mm,top=1.6mm,bottom=1.6mm,
  toptitle=1mm,bottomtitle=0.7mm,fonttitle=\bfseries,
  attach boxed title to top left={xshift=2.5mm,yshift=-1.15mm},
  boxed title style={colback=white,colframe=white,boxrule=0pt,left=1mm,right=1mm}}}
\newtcolorbox{zadaniBox}[1][]{commonlesson,title={Zadání},#1}
\newtcolorbox{prikladBox}[1][]{commonlesson,title={Příklad},#1}
\newtcolorbox{definiceBox}[1][]{commonlesson,title={Definice},#1}

% Spolecny box pro zadani a jeho reseni, bez mezery mezi castmi.
% Pouziti:
%   \begin{prikladReseniBox}
%     Zadani ulohy.
%     \tcblower
%     \begin{align*} ... \end{align*}
%   \end{prikladReseniBox}
% Spodni cast ma jemny sedy podklad a delici linku; lze v ni sazet
% vice radku matematiky. Obe casti maji jen jeden spolecny ramecek.
\tcbset{solutionlesson/.style={
  commonlesson,
  enhanced jigsaw,
  breakable,
  colbacklower=black!3,
  segmentation style={solid,draw=black!25,line width=0.35pt},
  fontlower=\small\color{notegray},
  before lower={\noindent\textbf{Řešení:}\enspace},
  leftlower=3mm,rightlower=3mm
}}
\newtcolorbox{prikladReseniBox}[1][]{solutionlesson,title={Příklad},#1}
\newtcolorbox{zadaniReseniBox}[1][]{solutionlesson,title={Zadání},#1}

% Pomucky pro autora pripravy: sedive poznamky; \reseni pro starsi pripravy
\newcommand{\poznamka}[1]{\par\begingroup\small\color{notegray}\itshape
  \textnormal{\textbf{Poznámka pro učitele.}}\enspace #1\par\endgroup}
\newcommand{\reseni}[1]{\par\begingroup\small\color{notegray}
  \textbf{Očekávané řešení:}\enspace #1\par\endgroup}
\newcommand{\blok}[1]{\needspace{5\baselineskip}\subsection*{#1}}
\newcommand{\vyrok}[1]{\uv{#1}}

\usepackage{tikz}
\usetikzlibrary{arrows.meta,positioning}
\nastavPripravu{1. ročník}{Funkce}{59}{Úvod do funkcí}

\begin{document}

\blok{Cíl a organizace hodiny}
Žáci si vytvoří přesnou představu funkce jako jednoznačného přiřazení. Naučí se číst základní zápisy a vypočítat funkční hodnotu i hledat argument z hodnoty. Seznámí se se čtyřmi způsoby zadání: tabulkou, předpisem, grafem a rekurencí. Důraz je na tom, že funkce nemusí být vzorec a že každému přípustnému vstupu musí odpovídat právě jeden výstup.

\poznamka{Doporučená skladba 45 minut je na konci přípravy. Zde uvedené oddíly označené „Jádro“ tvoří výběr pro danou hodinu; ostatní vzorové příklady, podrobnější vysvětlení a rezervy slouží učiteli podle tempa třídy.}

\blok{Motivace: vstup a výstup}
Každé kladné délce strany čtverce $a$ odpovídá právě jeden obsah $S=a^2$. Číslu $a=3$ odpovídá obsah $S=9$, číslu $a=4$ obsah $S=16$. Přiřazení má jasný vstup, výstup i pravidlo. Představme si naopak, že by jednomu vstupu $a=3$ odpovídaly dva různé obsahy. Pak by pravidlo neurčovalo výsledek jednoznačně.

\blok{1. Funkce jako přiřazení}
\subsection*{A. Kartézský součin a množina uspořádaných dvojic}
\begin{definiceBox}
Pro množiny $A$ a $B$ je jejich \textbf{kartézský součin} množina všech uspořádaných dvojic, jejichž první složka patří do $A$ a druhá do $B$:
\[
A\times B=\{(x,y)\mid x\in A,\ y\in B\}.
\]
Pořadí složek je důležité: obecně $(x,y)\ne(y,x)$.
\end{definiceBox}

\begin{definiceBox}
\textbf{Funkce} $f$ z množiny $A$ do množiny $B$ je podmnožina $A\times B$, pro kterou platí
\[
\forall x\in A\;\exists!y\in B:\ (x,y)\in f.
\]
Symbol $\forall$ znamená „pro každý“, $\exists$ „existuje“ a $\exists!$ „existuje právě jeden“. Každý prvek $A$ se tedy musí objevit jako první složka právě jedné dvojice; druhá složka této dvojice je jeho funkční hodnota.
\end{definiceBox}

\textbf{Konkrétní příklad.} Nechť $A=\{1,2\}$ a $B=\{a,b,c\}$. Potom
\[
A\times B=\{(1,a),(1,b),(1,c),(2,a),(2,b),(2,c)\}.
\]
Podmnožina $f=\{(1,a),(2,c)\}$ je funkcí: pro každý vstup z $A$ obsahuje právě jednu dvojici. Naproti tomu $R=\{(1,a),(1,b),(2,c)\}$ funkcí není, protože vstupu $1$ přiřazuje dvě různé hodnoty. Také $S=\{(1,a)\}$ není funkcí z $A$ do $B$: vstup $2$ nemá žádný výstup.

\subsection*{B. Funkce jako zobrazení}
\begin{definiceBox}
\textbf{Funkce} $f\colon A\to B$ je zobrazení, které každému prvku $x$ množiny $A$ přiřazuje právě jeden prvek $y$ množiny $B$. Píšeme $f(x)=y$ nebo $x\mapsto y$.
\end{definiceBox}

Tato definice říká totéž jako podmínka pro množinu dvojic: pro každý vstup existuje výstup a tento výstup je jediný. Množina $A$ je množina vstupů; nazývá se definiční obor. Množina $B$ je cílová množina. Skutečné výstupy mohou tvořit jen její část, tzv. obor hodnot; ten budeme přesně určovat v další hodině. Funkce nemusí pokrýt všechny prvky cílové množiny a různé vstupy mohou mít stejný výstup. Bijekce je zvláštní typ funkce s dalšími vlastnostmi, nikoli obecná definice funkce.

\begin{center}
\begin{tikzpicture}[x=1cm,y=1cm,>=Latex, every node/.style={font=\small}]
\draw (0,0) ellipse (0.72 and 1.12); \draw (3.15,0) ellipse (0.72 and 1.12);
\node at (0,1.45) {$A$}; \node at (3.15,1.45) {$B$};
\foreach \y/\lab in {0.65/1,0/2,-0.65/3} {\fill (0,\y) circle (1.5pt);\node[left=3pt] at (0,\y) {$\lab$};}
\foreach \y/\lab in {0.65/a,0/b,-0.65/c} {\fill (3.15,\y) circle (1.5pt);\node[right=3pt] at (3.15,\y) {$\lab$};}
\draw[->] (0.12,0.65)--(2.99,0.65); \draw[->] (0.12,0)--(2.99,-0.65); \draw[->] (0.12,-0.65)--(2.99,-0.65);
\node at (1.58,-1.55) {Funkce: každý vstup má právě jednu šipku.};
\end{tikzpicture}
\qquad
\begin{tikzpicture}[x=1cm,y=1cm,>=Latex, every node/.style={font=\small}]
\draw (0,0) ellipse (0.72 and 1.12); \draw (3.15,0) ellipse (0.72 and 1.12);
\node at (0,1.45) {$A$}; \node at (3.15,1.45) {$B$};
\foreach \y/\lab in {0.65/1,0/2,-0.65/3} {\fill (0,\y) circle (1.5pt);\node[left=3pt] at (0,\y) {$\lab$};}
\foreach \y/\lab in {0.65/a,0/b,-0.65/c} {\fill (3.15,\y) circle (1.5pt);\node[right=3pt] at (3.15,\y) {$\lab$};}
\draw[->] (0.12,0.65)--(2.99,0.65); \draw[->] (0.12,0)--(2.99,0.65); \draw[->] (0.12,-0.65)--(2.99,-0.65);
\node at (1.58,-1.55) {Dva různé vstupy mohou mít stejný výstup.};
\end{tikzpicture}
\end{center}

\begin{center}
\begin{tikzpicture}[x=1cm,y=1cm,>=Latex, every node/.style={font=\small}]
\draw (0,0) ellipse (0.72 and 1.12); \draw (3.15,0) ellipse (0.72 and 1.12);
\node at (0,1.45) {$A$}; \node at (3.15,1.45) {$B$};
\foreach \y/\lab in {0.65/1,0/2,-0.65/3} {\fill (0,\y) circle (1.5pt);\node[left=3pt] at (0,\y) {$\lab$};}
\foreach \y/\lab in {0.65/a,0/b,-0.65/c} {\fill (3.15,\y) circle (1.5pt);\node[right=3pt] at (3.15,\y) {$\lab$};}
\draw[->] (0.12,0.65)--(2.99,0.65); \draw[->,bend left=12] (0.12,0.08) to (2.99,0.65); \draw[->] (0.12,0)--(2.99,0); \draw[->] (0.12,-0.65)--(2.99,-0.65);
\node at (1.58,-1.55) {Není funkce: vstup $2$ má dva výstupy.};
\end{tikzpicture}
\qquad
\begin{tikzpicture}[x=1cm,y=1cm,>=Latex, every node/.style={font=\small}]
\draw (0,0) ellipse (0.72 and 1.12); \draw (3.15,0) ellipse (0.72 and 1.12);
\node at (0,1.45) {$A$}; \node at (3.15,1.45) {$B$};
\foreach \y/\lab in {0.65/1,0/2,-0.65/3} {\fill (0,\y) circle (1.5pt);\node[left=3pt] at (0,\y) {$\lab$};}
\foreach \y/\lab in {0.65/a,0/b,-0.65/c} {\fill (3.15,\y) circle (1.5pt);\node[right=3pt] at (3.15,\y) {$\lab$};}
\draw[->] (0.12,0.65)--(2.99,0.65); \draw[->] (0.12,0)--(2.99,-0.65);
\node at (1.58,-1.55) {Není funkce: vstup $3$ nemá výstup.};
\end{tikzpicture}
\end{center}

\subsection*{C. Slovní představa}
Funkce je pravidlo, které každému přípustnému vstupu přiřadí právě jeden výstup. Vstup vložíme do „přiřazovacího stroje“ a dostaneme jednoznačný výsledek. Množina dvojic zapisuje všechny vstupně-výstupní dvojice; zobrazení popisuje tutéž věc jako pravidlo přiřazení. Jednoznačnost neznamená, že výstupy musí být různé: dvěma vstupům může vyjít stejná hodnota.

\blok{2. Značení funkce}
\begin{center}
\small
\begin{tabularx}{\linewidth}{@{}p{35mm}X@{}}
\toprule
Zápis & Význam a použití\\
\midrule
$f$ & Označení funkce jako matematického objektu, tedy celého přiřazení.\\
$f(x)$ & Hodnota funkce $f$ při vstupu $x$; $x$ je argument. Nejde o součin $f\cdot x$.\\
$y=f(x)$ & Proměnná $x$ je nezávislá, proměnná $y$ závislá; rovnost říká, jak spolu souvisejí.\\
$f(x)=2x+3$ & Předpis: dosazením argumentu do pravé strany získáme funkční hodnotu.\\
$y=2x+3$ & Tentýž předpis vyjádřený pomocí proměnných $x,y$; body grafu splňují tuto rovnici.\\
$f\colon x\mapsto 2x+3$ & Zápis přiřazení: funkce $f$ přiřazuje vstupu $x$ hodnotu $2x+3$.\\
$f\colon A\to B$ & Určuje, že vstupy patří do $A$ a výstupy do cílové množiny $B$.\\
\bottomrule
\end{tabularx}
\end{center}

Je třeba rozlišit funkci $f$ (celé přiřazení), argument $x$ (vstup), hodnotu $f(x)$ (konkrétní výstup) a předpis, například $2x+3$ (způsob výpočtu hodnot). Předpis sám o sobě bez určení přípustných vstupů nemusí vždy určovat tutéž funkci; zatím však pracujeme s jednoduchými předpisy a jejich dosazováním.

\begin{prikladReseniBox}
Je dáno $f(x)=2x+3$. Vypočtěte $f(0)$, $f(2)$, $f(-2)$ a $f(-1)$. Poté určete argument, jehož funkční hodnota je $7$.
\tcblower
\begin{align*}
f(0)&=2\cdot0+3=3, & f(2)&=2\cdot2+3=7,\\
f(-2)&=2\cdot(-2)+3=-1, & f(-1)&=2\cdot(-1)+3=1.
\end{align*}
Hledáme $x$, pro které $f(x)=7$: řešíme $2x+3=7$, tedy $2x=4$ a $x=2$. Zpětná kontrola: $f(2)=7$.
\end{prikladReseniBox}

\begin{prikladReseniBox}
Pro $g(x)=x^2-1$ vypočtěte $g(-3)$ a určete všechny argumenty, pro které je $g(x)=8$.
\tcblower
\[
g(-3)=(-3)^2-1=9-1=8.
\]
Při hledání argumentu řešíme $x^2-1=8$, tedy $x^2=9$. Protože $3^2=9$ i $(-3)^2=9$, vycházejí dva argumenty: $x=3$ a $x=-3$. Jedna funkční hodnota tedy může pocházet z více vstupů.
\end{prikladReseniBox}

\poznamka{Při dosazování záporného čísla vždy nahrazujeme proměnnou celým číslem v závorce. Například $(-2)^2=4$, zatímco bez závorek by zápis $-2^2$ znamenal $-(2^2)=-4$.}

\blok{3. Způsoby zadání funkce}
\subsection*{A. Zadání tabulkou}
Tabulka udává vstupy a jejich odpovídající výstupy. Je-li řečeno, že definiční obor tvoří právě vypsané vstupy, úplná tabulka určuje celou funkci na tomto konečném oboru. Několik naměřených či vypočtených hodnot naproti tomu představuje jen známé dvojice; bez další informace z nich nelze odvodit jediný obecný předpis.

\begin{prikladReseniBox}
Funkce $t$ je zadána na $D(t)=\{-2,0,1,3\}$ úplnou tabulkou. Určete $t(1)$, argumenty s hodnotou $4$ a hodnotu $t(-2)$.
\[
\begin{array}{c|rrrr}
x&-2&0&1&3\\\hline
t(x)&4&1&4&-2
\end{array}
\]
\tcblower
Z tabulky přímo čteme $t(1)=4$ a $t(-2)=4$. Hodnotu $4$ mají argumenty $-2$ a $1$. Tabulka je úplná na zadané množině čtyř argumentů; neříká nic o hodnotě například v $x=2$, protože $2\notin D(t)$.
\end{prikladReseniBox}

\begin{prikladReseniBox}
Naměřené dvojice jsou $(0,1),(1,3),(2,5)$. Lze z nich bez dalšího předpokladu dokázat, že předpis je $f(x)=2x+1$?
\tcblower
Nelze. Známe jen tři naměřené hodnoty. Předpis $2x+1$ jim odpovídá, ale třeba při dalším měření v $x=3$ může vyjít jiná hodnota než $7$. Bez údaje o pravidlu tedy nevíme, co se děje mezi měřenými vstupy ani za nimi. Je-li zadáním funkce pouze na konečné množině $\{0,1,2\}$, tabulka ji na tomto oboru určuje úplně.
\end{prikladReseniBox}
\poznamka{Učitelská rezerva k přesnému zdůvodnění: také předpis $p(x)=2x+1+x(x-1)(x-2)$ dává pro $x=0,1,2$ hodnoty $1,3,5$, ale $p(3)=13$, zatímco $2\cdot3+1=7$.}

\subsection*{B. Zadání předpisem}
Předpis sděluje, jak vypočítat výstup ze vstupu. Jednoduché příklady jsou $f(x)=2x+1$, $g(x)=x^2$ nebo $h(x)=3-x$. Pro každý přípustný argument dosadíme za proměnnou; opačná úloha „pro který argument má funkce hodnotu $c$?“ vede k rovnici.

\begin{prikladReseniBox}
Je dáno $p(x)=3-2x$. Vypočtěte $p(0)$, $p(4)$ a $p(-2)$. Určete argument, pro který je $p(x)=9$.
\tcblower
\[
p(0)=3,\qquad p(4)=3-8=-5,\qquad p(-2)=3-2(-2)=7.
\]
Pro opačnou úlohu řešíme $3-2x=9$, a tedy $-2x=6$, $x=-3$. Kontrola: $p(-3)=3+6=9$.
\end{prikladReseniBox}

\begin{prikladReseniBox}
Je dáno $q(x)=(x-2)^2+1$. Vypočtěte $q(0)$ a určete argumenty, pro které $q(x)=10$.
\tcblower
\[
q(0)=(0-2)^2+1=5.
\]
Z rovnice $(x-2)^2+1=10$ dostáváme $(x-2)^2=9$, tedy $x-2=3$ nebo $x-2=-3$. Řešení jsou $x=5$ a $x=-1$. Kontrola: v obou případech je druhá mocnina rozdílu $9$, takže výstup je $10$.
\end{prikladReseniBox}

\begin{prikladReseniBox}[title={Volitelné rozšíření: náklady na výrobu}]
Výrobní náklady v korunách popisuje $C(n)=150n+400$, kde $n\in\mathbb N_0$ je počet výrobků. Co znamenají členy předpisu? Určete $C(0)$ a $C(5)$.
\tcblower
Argument $n$ je počet výrobků, funkční hodnota $C(n)$ jsou celkové náklady v Kč. Člen $400$ představuje fixní náklady i při nulové výrobě, člen $150n$ náklady na $n$ výrobků po $150$ Kč za kus. Proto $C(0)=150\cdot0+400=400$ Kč a $C(5)=150\cdot5+400=1150$ Kč. Počet výrobků je celé nezáporné číslo.
\end{prikladReseniBox}

\subsection*{C. Zadání grafem}
Graf funkce je množina bodů $(x,f(x))$ v souřadnicové rovině. První souřadnice bodu je argument, druhá jeho funkční hodnota. K určení $f(a)$ vedeme od $x=a$ svisle ke grafu a od průsečíku vodorovně k ose $y$.

\begin{center}
\begin{tikzpicture}[x=0.75cm,y=0.55cm,>=Latex]
\draw[->] (-3.5,0)--(3.8,0) node[right] {$x$};
\draw[->] (0,-3.5)--(0,4.2) node[above] {$y$};
\draw[step=1,gray!25,very thin] (-3,-3) grid (3,4);
\foreach \x in {-3,-2,-1,1,2,3}{\draw (\x,0.08)--(\x,-0.08) node[below=2pt] {\small $\x$};}
\foreach \y in {-3,-2,-1,1,2,3,4}{\draw (0.08,\y)--(-0.08,\y) node[left=2pt] {\small $\y$};}
\draw[thick,black,domain=-3:3,samples=2] plot (\x,{\x+1});
\fill ( -2,-1) circle (2.4pt) node[below left=3pt] {\small $(-2;-1)$};
\fill (0,1) circle (2.4pt) node[above right=3pt] {\small $(0;1)$};
\fill (2,3) circle (2.4pt) node[above left=3pt] {\small $(2;3)$};
\end{tikzpicture}
\end{center}

Na přímce $y=x+1$ jsou například body $(-2,-1)$, $(0,1)$ a $(2,3)$, takže $f(-2)=-1$, $f(0)=1$ a $f(2)=3$. Svislicový test vychází z definice: svislice $x=a$ zachycuje body se stejným argumentem $a$. Protne-li obrázek ve dvou různých bodech, mají tyto body stejnou první souřadnici, ale různé druhé souřadnice; jednomu argumentu by tak odpovídaly dvě hodnoty. Každá svislice proto smí graf funkce protínat nejvýše jednou.

\begin{center}
\begin{tikzpicture}[x=0.75cm,y=0.75cm,>=Latex]
\draw[->] (-2.8,0)--(3.2,0) node[right] {$x$};
\draw[->] (0,-2.8)--(0,3.2) node[above] {$y$};
\draw[thick,black] (0,0) circle (2);
\draw[dashed,thick] (1.2,-2.65)--(1.2,2.65) node[above] {$x=1{,}2$};
\fill (1.2,1.6) circle (2.5pt) node[above right=3pt] {\small $(1{,}2;1{,}6)$};
\fill (1.2,-1.6) circle (2.5pt) node[below right=3pt] {\small $(1{,}2;-1{,}6)$};
\end{tikzpicture}
\end{center}

Kružnice $x^2+y^2=4$ svislici $x=1{,}2$ protne ve dvou bodech, protože $y=\pm\sqrt{4-1{,}44}=\pm1{,}6$. Není tedy grafem funkce $y$ závislé na $x$ na celém zobrazeném oboru. Jde o závislost, v níž některému vstupu odpovídají dvě hodnoty.

\begin{prikladReseniBox}
Z grafu přímky $y=x+1$ určete $f(-1)$ a argument, pro který je funkční hodnota $3$.
\tcblower
Na přímce při $x=-1$ vychází $y=-1+1=0$, tedy $f(-1)=0$. Pro hodnotu $3$ řešíme $x+1=3$, takže $x=2$; v grafu jde o bod $(2,3)$.
\end{prikladReseniBox}

\subsection*{D. Zadání rekurentně}
Rekurentní předpis určuje další hodnotu pomocí jedné či více předchozích. Musí obsahovat počáteční hodnotu; bez ní nelze určit, odkud výpočet začne. Posloupnost je funkce, jejíž vstupy tvoří přirozená čísla od zadaného počátku, tedy diskrétní obor.

\begin{prikladReseniBox}
Je dáno $a_1=2$ a $a_{n+1}=a_n+3$ pro $n\geq1$. Vypočtěte prvních pět členů.
\tcblower
\begin{align*}
a_1&=2,\\
a_2&=a_1+3=2+3=5,\\
a_3&=a_2+3=5+3=8,\\
a_4&=a_3+3=8+3=11,\\
a_5&=a_4+3=11+3=14.
\end{align*}
Každý člen kromě prvního získáme přičtením $3$ k předchozímu: $2,5,8,11,14,\ldots$.
\end{prikladReseniBox}

\begin{prikladReseniBox}
Posloupnost je zadána $b_1=3$, $b_{n+1}=2b_n$. Určete první čtyři členy a vysvětlete, co by chybělo při vynechání údaje $b_1=3$.
\tcblower
\[
b_1=3,\quad b_2=2\cdot3=6,\quad b_3=2\cdot6=12,\quad b_4=2\cdot12=24.
\]
Bez počáteční hodnoty by pravidlo určovalo pouze vztah mezi sousedními členy, nikoli konkrétní posloupnost; například počáteční hodnoty $1$ a $3$ vedou k různým posloupnostem.
\end{prikladReseniBox}

\blok{4. Procvičení v několika úrovních}
\subsection*{Jádro: společné procvičení}
\begin{zadaniReseniBox}
\begin{enumerate}[label=\alph*)]
\item Rozhodněte, zda $R=\{(1,a),(2,b),(2,c),(3,a)\}$ určuje funkci z $\{1,2,3\}$ do $\{a,b,c\}$.
\item Pro $u(x)=x^2+2$ určete $u(-2)$ a všechny argumenty s funkční hodnotou $6$.
\item Úplná tabulka na oboru $\{-1,0,2\}$ má hodnoty $v(-1)=4$, $v(0)=1$, $v(2)=4$. Kolik argumentů má hodnotu $4$? Je tím určena $v(1)$?
\item Rekurentně je dáno $c_1=1$, $c_{n+1}=c_n+4$. Určete $c_2,c_3,c_4$.
\end{enumerate}
\tcblower
\begin{enumerate}[label=\alph*)]
\item Ne. Vstupu $2$ odpovídají dvě různé hodnoty $b$ a $c$.
\item $u(-2)=(-2)^2+2=6$. Rovnice $x^2+2=6$ dává $x^2=4$, tedy $x=-2$ nebo $x=2$.
\item Dva argumenty, $-1$ a $2$. Hodnota $v(1)$ určena není, neboť $1$ neleží v uvedeném definičním oboru.
\item $c_2=1+4=5$, $c_3=5+4=9$, $c_4=9+4=13$.
\end{enumerate}
\end{zadaniReseniBox}

\subsection*{Rozšíření: argumentace a propojení reprezentací}
\begin{zadaniReseniBox}
\begin{enumerate}[label=\alph*)]
\item Uvažujte funkci $w\colon\{1,2,3\}\to\{p,q,r,s\}$ danou dvojicemi $\{(1,p),(2,p),(3,r)\}$. Je to funkce? Jaká je cílová množina a které její prvky jsou skutečnými výstupy?
\item Tabulka udává $r(-1)=3$, $r(0)=1$, $r(1)=-1$. Odpovídají hodnoty předpisu $r(x)=1-2x$? Co lze z tabulky říci o případných dalších argumentech?
\item Je množina bodů $\{(0,1),(1,2),(2,3),(1,4)\}$ grafem funkce proměnné $x$? Zdůvodněte bez kreslení.
\item Předpis $m(x)=5-2x$. Určete argument s hodnotou $1$ a ověřte výsledek dosazením.
\end{enumerate}
\tcblower
\begin{enumerate}[label=\alph*)]
\item Ano, každý ze tří vstupů má právě jednu hodnotu. Cílová množina je $\{p,q,r,s\}$; mezi skutečné výstupy patří jen $p$ a $r$. Prvek $q$ ani $s$ funkce nepřiřazuje žádnému vstupu.
\item Ano: $1-2(-1)=3$, $1-2\cdot0=1$ a $1-2\cdot1=-1$. Tabulka ale sama nedokazuje, že tento předpis platí také mimo uvedené tři argumenty.
\item Ne, protože argumentu $1$ odpovídají dvě hodnoty, $2$ a $4$; svislice $x=1$ by protnula množinu ve dvou bodech.
\item Řešíme $5-2x=1$, takže $x=2$. Kontrola: $m(2)=5-4=1$.
\end{enumerate}
\end{zadaniReseniBox}

\subsection*{Rezervní úlohy}
\begin{zadaniReseniBox}
\begin{enumerate}[label=\alph*)]
\item Je $P=\{(a,1),(b,1),(c,2)\}$ funkcí z $\{a,b,c\}$ do $\{1,2,3\}$? Vysvětlete rozdíl mezi cílovou množinou a hodnotami, které se skutečně objeví.
\item Pro $z(x)=2(x-3)+1$ vypočtěte $z(0)$ a určete argument s hodnotou $-5$.
\item Kružnice $x^2+y^2=9$ je při pohledu na $y$ jako funkci proměnné $x$ testována svislicí $x=0$. Kolik průsečíků vznikne a jaké jsou odpovídající hodnoty $y$?
\item Je posloupnost $d_1=5$, $d_{n+1}=d_n-2$ jednoznačně určena? Vypočtěte pět prvních členů.
\end{enumerate}
\tcblower
\begin{enumerate}[label=\alph*)]
\item Ano. Každý vstup $a,b,c$ má jedinou hodnotu. Cílová množina je $\{1,2,3\}$, skutečné výstupy jsou $\{1,2\}$; hodnota $3$ se nevyskytuje.
\item $z(0)=2(0-3)+1=-6+1=-5$. Z rovnice $2(x-3)+1=-5$ plyne $2(x-3)=-6$, $x-3=-3$, a tedy $x=0$; kontrola souhlasí.
\item Dva průsečíky: z $0+y^2=9$ dostaneme $y=3$ a $y=-3$. Kružnice proto není grafem jedné funkce $y(x)$ na celém oboru.
\item Ano, počáteční hodnota i pravidlo jsou zadány. Členy jsou $5,3,1,-1,-3$.
\end{enumerate}
\end{zadaniReseniBox}

\blok{5. Domácí práce}
\begin{zadaniReseniBox}
\begin{enumerate}[label=\alph*)]
\item Pro $h(x)=3x-2$ vypočtěte $h(0)$, $h(-2)$ a určete argument, pro který $h(x)=10$.
\item Funkce $k$ je na oboru $\{0,1,2,3\}$ zadána tabulkou $k(0)=2$, $k(1)=0$, $k(2)=2$, $k(3)=-1$. Vypište dvojice, které tvoří funkci, a určete argumenty s hodnotou $2$.
\item Rozhodněte, zda množina $\{(-1,2),(0,1),(1,2),(1,3)\}$ může být grafem funkce. Uveďte důvod.
\item Rekurentně je dáno $e_1=4$, $e_{n+1}=e_n+5$. Určete prvních pět členů.
\end{enumerate}
\tcblower
\begin{enumerate}[label=\alph*)]
\item $h(0)=-2$, $h(-2)=3(-2)-2=-8$. Z $3x-2=10$ plyne $3x=12$, tedy $x=4$; kontrola $h(4)=10$.
\item Dvojice jsou $\{(0,2),(1,0),(2,2),(3,-1)\}$. Hodnotu $2$ mají argumenty $0$ a $2$.
\item Ne. Argumentu $1$ jsou přiřazeny dvě různé hodnoty, $2$ a $3$.
\item $e_1=4$, $e_2=9$, $e_3=14$, $e_4=19$, $e_5=24$.
\end{enumerate}
\end{zadaniReseniBox}

\blok{6. Přechod k vlastnostem funkcí}
V následujících hodinách budeme funkce popisovat podrobněji. Budeme určovat \textbf{definiční obor} (pro které vstupy je funkce zadána) a \textbf{obor hodnot} (kterých výstupů skutečně nabýváme); sledovat \textbf{monotonii} (zda funkce při růstu argumentu roste či klesá); zkoumat \textbf{sudost a lichost} a \textbf{prostost}; hledat \textbf{maximum a minimum}; a rozhodovat, zda jsou hodnoty \textbf{omezené shora nebo zdola}. V této hodině tyto pojmy pouze pojmenujeme. Grafy a předpisy nám budou poskytovat různé pohledy na tentýž objekt.

\poznamka{Časový odhad 45minutového jádra včetně organizace: motivace 3 min; kartézský součin a definice s jedním diagramem 10 min; slovní představa a notace 8 min; předpis a dva výpočty 8 min; ukázková tabulka a čtení jediného bodu z grafu 7 min; krátké procvičení s kontrolou 7 min; shrnutí a zadání domácí práce 2 min. Součet je 45 min. Zbylé diagramy, druhý graf, rekurence, rozšíření a rezervy jsou volitelné nebo určené k navázání. Pokud učitel chce v této hodině důkladně předvést všechny čtyři způsoby zadání a řešit více příkladů, celý dokument přesáhne jednu vyučovací hodinu; není vhodné urychlovat jádro, proto volitelný obsah vybírej podle třídy. Domácí práce je přibližně na 15–20 min.}

\end{document}
